% Aporte de Darwin Joel Sánchez Vera.
\subsection{Discusión de latencia y rendimiento}

TCP obtuvo menor latencia en las cuatro condiciones. La diferencia más grande apareció con mensajes de 100~B en modo secuencial: la mediana fue de 0,0953~ms con TCP y de 0,6171~ms con gRPC, es decir, gRPC tuvo una mediana 547,6~\% mayor. En modo concurrente con 10~KiB la diferencia de la mediana bajó a 27,9~\%, aunque el percentil 95 de gRPC siguió siendo 145,4~\% mayor que el de TCP. En la figura~\ref{fig:barras} los intervalos de la latencia media de TCP y gRPC no se superponen en tres condiciones. Solo se superponen con 10~KiB en modo concurrente, por lo que en esa condición la diferencia de medias es menos concluyente que la de medianas.

El tamaño del mensaje afectó de forma distinta a cada mecanismo. En modo secuencial, al pasar de 100~B a 10~KiB la mediana de TCP subió de 0,0953 a 0,1869~ms (96~\% más), mientras que la de gRPC pasó de 0,6171 a 0,6054~ms. Con nuestras implementaciones, el costo de gRPC parece depender poco de la carga en ese rango y mucho del trabajo fijo de cada llamada; en TCP, en cambio, codificar y copiar el JSON crece con el mensaje. No medimos por separado la serialización, así que esta explicación queda como hipótesis a comprobar.

La concurrencia también tuvo efectos opuestos. En TCP, enviar desde los tres nodos a la vez elevó la mediana con 100~B de 0,0953 a 0,3014~ms. Los tres nodos de TCP comparten un solo proceso, lo que puede explicar parte de ese aumento. En gRPC la mediana concurrente fue menor que la secuencial (0,5166 frente a 0,6171~ms con 100~B); nuestras mediciones no permiten aislar la causa de ese comportamiento.

El rendimiento de gRPC fue entre 55,9~\% y 89,6~\% menor que el de TCP. Esta comparación no es tan directa como la del RTT. En gRPC, la duración de la tanda incluye, por cada envío, la llamada del lanzador al servicio \texttt{ControlNodo.Intercambiar}, que no forma parte del RTT medido. Con 100~B en modo secuencial, cada envío de gRPC ocupó en promedio 1,21~ms de la tanda ($1/825{,}6$~s), aunque su RTT medio fue de 0,655~ms; en TCP la diferencia equivalente fue de 0,126 frente a 0,115~ms. Esto es coherente con que el rendimiento de gRPC suba en modo concurrente (2\,112,9 frente a 825,6~mensajes/s), donde las tres llamadas de control se solapan.

La variación entre corridas fue mayor en gRPC con 10~KiB. Las corridas de la repetición 2 de esas dos combinaciones, ejecutadas una tras otra en las posiciones 22 y 23 de la sesión, tuvieron medias de 1,49 y 1,77~ms, frente a valores entre 0,48 y 0,83~ms en las demás repeticiones. Esas dos corridas amplían el intervalo de confianza de la figura~\ref{fig:barras}. Las conservamos porque son mediciones válidas y el protocolo no prevé descartar corridas por su valor.

En la hipótesis planteamos que TCP tendría menor mediana con 100~B y que con 10~KiB la serialización podría cambiar la diferencia. La primera parte se cumplió. La segunda se cumplió parcialmente: con 10~KiB la diferencia relativa de la mediana bajó de 547,6~\% a 224,0~\% en modo secuencial y de 71,4~\% a 27,9~\% en modo concurrente, pero TCP mantuvo la menor mediana en todas las condiciones.
